\begin{abstract}
A current direction influences a material inverse update through material excitation, multiple-scattering feedback, and reception. For a specified Galerkin current-space change, we expose the retained feedback and factor the Jacobian change into observation and injection signatures. Their paired material pullbacks contain the complete regularized Gauss--Newton step change. The change acts on both data curvature and the residual-dependent normal-equation term, so increased information volume need not improve the step. This mechanism leads to receiver-centered conditional exchange: retain a receiver-derived backbone, rank actual model-induced displacements using a richer reduced reference, and check finalists on the complete frozen quadratic using cached output and directional tangent actions. In three-dimensional vector-Maxwell tests on four geometries with Gaussian and voxel material spaces, at most three exchanges improve 35 of 36 state--budget cases and leave one unchanged. The median object-summed quadratic-gap ratio is 0.348, a 65.2\% reduction at unchanged current dimension. Aggregating one-decision gains at each object's common checkpoints, full-dictionary proxy ranking captures 98.1--99.6\% of the best local one-decision gain; restricted candidate coverage explains the largest missed opportunities. Maxwell witnesses also show increasing information volume with negative step gain and positive joint gain from two unfavorable singleton additions. Complete costs and earlier Krylov controls locate the practical limits: these results establish a conditional current-design method for accurate frozen material updates, with candidate acquisition, output-error certification, and nonlinear transfer determining its deployment value.
\end{abstract}
